\typeout{NT FILE chapter8.tex}%

\chapter{OFLAT Implementation}
\label{cha:oflat-implementation}

This chapter describes how finite-state transducer (FST) functionality is integrated into the existing OFLAT code. It
covers module responsibilities, rendering logic, and how the new pieces reuse pre-existing abstractions for
other automata models. The goal is to provide FST support without altering the architecture established for FA, PDA,
and TM. The design perspective is covered separately in \autoref{cha:oflat-design}.

\section{Workspace and Graph Interaction}
\label{sec:oflat_workspace}
OFLAT keeps a shared rendering layer around Cytoscape.js, with thin OCaml bindings in \texttt{Cytoscape.ml}. FST support
reuses the same primitives as finite automata:
\begin{itemize}
  \item \textbf{Node lifecycle:} methods such as \texttt{addNode}, \texttt{addInitialNode}, and \texttt{addFinalNode}
      extend \texttt{TransducerView.model} (which inherits \texttt{AutomatonView.model}). They update the in-memory
      representation and immediately reflect changes on the graph (initial/accepting glyphs, labels).
  \item \textbf{Edges with input and output:} \texttt{inputEdges} maps each transition \((src, inSym, outSym, dst)\) to a
      Cytoscape edge labeled \texttt{in:out}. Epsilon is rendered with the same tilde used elsewhere in OFLAT to stay
      consistent with existing models.
  \item \textbf{Style reset and repaint:} before runs, styles are reset (\texttt{Cytoscape.faStyle}) so later execution
      highlights remain visible. Minimization coloring and best-path accents reuse the same paint helpers used for FA.
\end{itemize}
\begin{figure}[H]
  \centering
  \includegraphics[width=0.95\textwidth]{fst_rendermin.png}
  \caption{Rendering view with input:output edges and minimized coloring.}
  \label{fig:oflat-impl-render}
\end{figure}

Edge construction for input/output labels:
\begin{verbatim}
method inputEdges (cy: Cytoscape.cytoscape Js_of_ocaml.Js.t) =
	let mapToCytoscapeEdge transitions =
		Set.map (fun (src, iSym, oSym, dst) ->
			let label = (symb2str iSym) ^ ":" ^ (symb2str oSym) in
			(src, label, dst)
		) self#representation.transitions
	in
	Set.iter (Cytoscape.addEdge cy) (mapToCytoscapeEdge self#representation.transitions)
\end{verbatim}
Each transition becomes a Cytoscape edge labeled \texttt{input:output}, using the same epsilon glyph as elsewhere.

\section{Information Box, Poppers, and Menus}
\label{sec:oflat_infobox}
FST execution introduces richer configuration data (state, remaining input, accumulated output). The UI represents this via:
\begin{itemize}
  \item \textbf{Configuration menus:} built with \texttt{contextMenus}, populated from configuration sets at each step.
	  IDs are generated per-state to allow selective show/hide on hover. Menus are rebuilt whenever the step changes.
  \item \textbf{Poppers for configuration counts:} \texttt{buildPoppersConfigsCounter} creates poppers per node showing
	  how many configurations currently live there. They track node movement and viewport changes via subscribed events
	  (position, pan, zoom, resize).
  \item \textbf{Info display:} configuration details (state, remaining input, output so far) are presented in the info box
	  so users can correlate graph highlights with exact words being processed.
\end{itemize}

\section{Step-by-Step Execution Pipeline}
\label{sec:oflat_stepbystep}
Execution leverages the library \texttt{Transducer.acceptFull}, which returns (accepted flag, best path, trail of
configuration sets). The view code keeps these in mutable fields and drives the UI as follows:
\begin{itemize}
  \item \textbf{Trail and path storage:} \texttt{steps} holds the trail (array of configuration sets); \texttt{bestPath}
	  stores a single accepting path. Statistics (visited configs, timing, exactness) come from \texttt{Model.stats()}.
  \item \textbf{Painting states:} each step paints all active states; accepting states use the same accent as FA. The
	  chosen accepting path (if any) is overpainted in a darker color so users see which branch was selected among
	  nondeterministic options.
  \item \textbf{Menus and poppers refresh:} on every step change, configuration menus are rebuilt and poppers refreshed to
	  reflect the current configuration counts, ensuring hover interactions stay accurate.
  \item \textbf{Word and output view:} the sentence display marks consumed vs. remaining input with a bar and shows the
	  accumulated output for the current configuration, aligning text feedback with graph highlights.
\end{itemize}

Initial and next steps:
\begin{verbatim}
method setInitialStep cy =
    self#staticAcceptFull;
    Cytoscape.resetStyle cy Cytoscape.faStyle;
    if Array.length steps > 0 then self#initAllMenusAndFeatures cy steps.(0)

method setNextStep cy =
    if position < Array.length steps then
        self#updateAllMenusAndFeatures cy steps.(position)
\end{verbatim}
Step initialization runs the computation, resets styles, and builds menus/poppers; 
advancing a step refreshes the overlays using cached trail data.

\subsection*{UI data flow and state}
The UI keeps a set of mutable fields in the model instance to avoid recomputation: (1) \texttt{steps} (the trail) and
	\texttt{bestPath}, (2) the current step index, (3) cached popper handles, and (4) the last rendered input/output strings.
On a step change the view invalidates only what changed: it repaints nodes/edges, rebuilds menus and poppers from the cached
trail slice, and reuses the existing Cytoscape elements instead of rebuilding the graph. Rendering throttles to the Cytoscape
animation frame callback to avoid flicker on rapid scrubbing.

\section{Tables, Traces, and Minimization}
\label{sec:oflat_tables}
Tabular views rely on the shared \texttt{HTMLTable} helpers and adapt to FST specifics:
\begin{itemize}
  \item \textbf{Transition tables:} built per input alphabet (plus epsilon if present). Cells aggregate destinations and
	  output symbols as \texttt{dest/output}. Rebuilds occur after edits to keep data consistent with the graph.
  \item \textbf{Execution traces:} for a best path, the trace table lists step index, current state, produced output, and
	  remaining input. The final row is colored green/red depending on acceptance.
	\item \textbf{Minimization coloring:} equivalence classes from the library minimization partition are painted on the
	  graph using the same color palette as FA; edges are redrawn using \texttt{inputEdges} to keep labels intact.
\end{itemize}

\section{Exercise Handling}
\label{sec:oflat_exercises}
Exercise evaluation for FSTs follows the same pipeline used for other models but is adapted to the transducer semantics:
\begin{itemize}
	\item \textbf{Input\,$\to$\,Output assertions:} Exercise words can be written as \texttt{input->output}. The view passes
		these strings to the library helper \texttt{acceptExpectedOutput}, which parses the arrow, runs the transducer on the
		input, and asserts the produced output matches exactly (and belongs to the declared output alphabet). Plain words
		without an arrow are handled by \texttt{acceptCheckOutput}, which enforces input and output alphabet membership.
	\item \textbf{Inside/Outside sets:} The inside set lists pairs the model must accept (e.g., \texttt{a->a}, \texttt{ab->ab});
		the outside set can include counterexamples (e.g., \texttt{a->b}) to ensure incorrect mappings are rejected. The UI
		surfaces failures by highlighting the offending cases and showing the actual output returned by the transducer.
	\item \textbf{Properties:} Standard properties (deterministic, complete, Mealy/Moore, minimal, clean) are delegated to the
		library checks already described. Results are shown alongside input/output verdicts so users see structural issues and
		behavioural mismatches together.
	\item \textbf{Feedback loop:} On submission, the exercise runner collects acceptance results, output mismatches, and property
		violations, then updates the info box and tables. Because the same acceptance helpers drive both the exercise evaluator
		and the step-by-step view, traces shown in the UI match exactly what the checker has used.
	\item \textbf{UI specifics:} Failing inside cases are shown with the actual output alongside the expected arrow form; failing
		outside cases are flagged when the model accepts or emits output where it should not. Property failures appear in the
		same panel so users see behavioural and structural issues together. The editor keeps the graph intact—only tables and
		status banners refresh—to reduce disruption while iterating on fixes.
\end{itemize}
\begin{figure}[H]
  \centering
  \includegraphics[width=0.95\textwidth]{ex_feedback.png} 
  \caption{Exercise feedback: input$\to$output cases with actual outputs and property results.}
  \label{fig:oflat-impl-exercises}
\end{figure}

Exercise hooks exposed by the transducer module:
\begin{verbatim}
let checkExercise ex fst = Model.checkExercise ex 
(acceptExpectedOutput fst) (checkProperty fst)
let checkExerciseFailures ex fst = Model.checkExerciseFailures 
ex (acceptExpectedOutput fst) (checkProperty fst)
\end{verbatim}
Exercises reuse the library's output-aware acceptance and property checks directly, so the UI surfaces the same 
verdicts as the core model.

\section{Integration Boundaries}
\label{sec:oflat_integration}
The FST view layer is intentionally thin: it delegates acceptance, generation, determinization, and minimization to library
modules (\texttt{Transducer.accept}, \texttt{acceptFull}, \texttt{generate}, \texttt{minimize}), then renders the returned
trail/partitions. This preserves OFLAT's existing separation between computation (library) and interaction (view).
We also leverage the newer exercise-facing helpers \texttt{acceptCheckOutput} and \texttt{acceptExpectedOutput} to keep
input/output validation consistent with the transducer library, ensuring UI traces and checks match the library's
acceptance semantics. When determinization fails to terminate (e.g., output conflicts causing an infinite construction),
the UI reports a non-termination message and leaves the existing graph intact rather than spinning indefinitely.

